\par\noindent\textbf{(1)} 机械角动量 \(\boldsymbol L=m\boldsymbol r\times\dot{\boldsymbol r}\)，其导数为 \(d\boldsymbol L/dt=\boldsymbol r\times(q\boldsymbol v\times\boldsymbol B)\)，球面法向约束不产生关于球心的力矩。用 \(\boldsymbol B=g_m\boldsymbol r/R^3\)、\(\boldsymbol r\cdot\boldsymbol v=0\) 及三重叉乘得 \(\dot{\boldsymbol L}=qg_m\boldsymbol v/R=qg_m\,d\hat{\boldsymbol r}/dt\)。
\par\noindent\textbf{(2)} 因此定义 \(\boldsymbol J=\boldsymbol L-qg_m\hat{\boldsymbol r}\)，直接得到 \(d\boldsymbol J/dt=0\)。由于 \(\boldsymbol L\perp\hat{\boldsymbol r}\)，有 \(J^2=L^2+(qg_m)^2\)，同时 \(\boldsymbol J\cdot\hat{\boldsymbol r}=-qg_m\)。这是一种含磁场角动量项的守恒量，不等于普通机械角动量本身。
\par\noindent\textbf{(3)} 守恒标量 \(\boldsymbol J\cdot\hat{\boldsymbol r}=-qg_m\) 意味着粒子位置矢量始终与固定 \(\boldsymbol J\) 形成固定夹角 \(\cos\theta=-qg_m/J\)。于是球面轨迹位于球 \(r=R\) 和以 \(\boldsymbol J\) 为轴的圆锥相交的圆上，只有 \(J>|qg_m|\) 时轨道为非退化圆周。
\par\noindent\textbf{(4)} 以 \(\boldsymbol J\) 为极轴，轴向分量满足 \(J=mR^2\sin^2\theta\,\dot\phi-qg_m\cos\theta\)。代入 \(\cos\theta=-qg_m/J\) 得 \(\dot\phi=J/(mR^2)\)，为均匀进动。由于 \(\boldsymbol v\cdot(q\boldsymbol v\times\boldsymbol B)=0\)、法向约束不做功，\(E=mv^2/2\) 守恒，与恒定角速率及定角轨迹一致。
